\paragraph{What the exercises do not show.}
Three limits qualify the results. The first is that $\bar R=b$ requires the market to be outside Complete Adverse Selection, which holds at the calibrated hiding cost by construction (Section~\ref{ss:threshold}). The second is that the zero effect of salience on renovation is an implication of one hiding cost per market. With costs that vary across sellers, an informed seller who strictly prefers hiding would lose her rent by renovating and disclosing, so making hiding less attractive could change her decision; we do not model that margin, and Section~\ref{ss:att_saving} bounds what it could change. The third is that the benchmark renovation costs are illustrative. The close of the exercise is the contrast the model delivers: salience changes who hides and, to first order, not who renovates, whereas a subsidy changes who renovates and not who hides, until the stock of disclosed poor ratings is exhausted (Proposition~\ref{prop:subsidy_ge}).
